\documentclass{ltjsarticle}

\usepackage{style/main}
\usepackage{fancyhdr}
\pagestyle{fancy}
\fancyhf{} 
\renewcommand{\headrulewidth}{0pt}
\renewcommand{\footrulewidth}{0pt}


\begin{document}

\recipientinfo
{各応援団代表者様}
{令和8年9月1日}
{高校放送部}

\doctitle{【重要】体育祭応援合戦 音源提出について}

\bodytext

体育祭応援合戦の実施につきまして、下記の通り使用音源の提出をお願いいたします。

\kicenter

\begin{enumerate}

\item 提出物

\begin{enumerate}
\item 音源を収録したCD 2枚（予備用・本番用）
\item 音源の楽曲情報・クレジットを記載した用紙
\end{enumerate}

\item 提出期限

令和8年9月11日（金）

\item 提出先

高校放送部顧問　谷口　または

高校放送部　梶原（高3－2）

\item 取り扱いと注意

\begin{enumerate}
\item CDはCD-DA規格で記録してください。
\item CDケースおよびCD本体に団体名と「本番」「予備」を記載してください。
\item 本番用CDは確認後ご返却します。本番当日、再度ご提出ください。
\item 予備用CDは本番まで放送部にてお預かりします。
\item 1トラックに1曲を収録し、各トラック先頭に3秒間の無音を挿入してください。
\item 本番用・予備用CDは、ともに同一内容（音源・構成）としてください。
\end{enumerate}

\end{enumerate}

\docend

お問い合わせ
\begin{itemize}
\item 高校放送部　体育祭音響統括責任者　梶原（高3－2）
\item 高校放送部　体育祭渉外班長　川合（高3－5）
\end{itemize}

\end{document}